\section{The canonical derived dual and the retained original tau complex}
\label{sec:canonical-coherent-bidual}

This continuation uses the actual sheaf morphism (OCQ.2), constructed
from the original theta quotient in A1427, and the original four-point
amplitude complex in A1523.  All completed objects below occur as
explicit targets in a derived calculation.  The original analytic
sheaf and the original Schwartz spaces are unchanged.

\subsection{The global splitting is an actual sheaf map}

Keep $X=\mathbb C$ with coordinate $s$, $g=2\xi$, and
\[
 0\longrightarrow\mathcal K\xrightarrow{\iota}\mathcal M
 \xrightarrow{\beta}\mathcal N\longrightarrow0,
 \qquad \mathcal N=\mathcal O_X/(g).
 \tag{DC.1}
\]
These are precisely the sheaves of (OCQ.1)--(OCQ.3).  At every
analytic stalk the original calculation makes every nonzero germ
invertible on $\mathcal K_{s_0}$.  In particular multiplication by
the same entire function $g$ is an isomorphism of the sheaf
$\mathcal K$: a morphism of sheaves is an isomorphism when its maps
on all stalks are isomorphisms.  The inverse is unique and commutes
with restrictions.

Since $g\mathcal N=0$, multiplication by $g$ on $\mathcal M$ factors
uniquely through a sheaf map $u_g:\mathcal M\to\mathcal K$ with
$\iota u_g=g\operatorname{id}_{\mathcal M}$.  Define
\[
 r=(g|_{\mathcal K})^{-1}u_g:\mathcal M\longrightarrow\mathcal K,
 \qquad P=\operatorname{id}_{\mathcal M}-\iota r.
 \tag{DC.2}
\]
One has $r\iota=\operatorname{id}_{\mathcal K}$, $rP=0$,
$P^2=P$, $\beta P=\beta$, and $gP=0$, by substituting (DC.2).
The intersection $\mathcal M[g]\cap\mathcal K$ is zero, because
$g$ acts injectively on $\mathcal K$.  Thus $\beta$ restricts to an
isomorphism $\mathcal M[g]\to\mathcal N$.  Its inverse
$j:\mathcal N\to\mathcal M$ satisfies $j\beta=P$, and gives the
canonical sheaf isomorphisms
\[
 (\beta,r):\mathcal M\xrightarrow{\sim}\mathcal N\oplus\mathcal K,
 \qquad (n,k)\longmapsto jn+\iota k.
 \tag{DC.3}
\]
Both composites are the identity by the displayed projection identities.
No local vector-space basis or choices of representatives enter these
maps.  In particular the stalk decomposition (OCQ.7) genuinely glues.

Let $\omega_X=\mathcal O_X\,ds$.  Applying the derived internal Hom
functor to the finite direct sum (DC.3) gives
\[
 R\mathcal H om(\mathcal M,\omega_X)
 \simeq R\mathcal H om(\mathcal N,\omega_X)
       \oplus R\mathcal H om(\mathcal K,\omega_X).
 \tag{DC.4}
\]
The original exact locally free resolution
$[\mathcal O_X\xrightarrow{g}\mathcal O_X]\to\mathcal N$,
with its free terms in degrees $-1,0$, computes the first summand.
Multiplication by $g$ is injective on $\omega_X$, so
\[
 R\mathcal H om(\mathcal N,\omega_X)
       \simeq(\omega_X/g\omega_X)[-1].
 \tag{DC.5}
\]
The injection of (DC.5) into (DC.4) is $\beta^*$ and its retraction
is $j^*$; their composite is the identity since $\beta j=1$.
This is a global derived summand of the actual source.  It makes no
identification between stalks of the second internal Hom in (DC.4)
and Ext groups computed from its noncoherent stalks.

\subsection{A canonical local dual without a choice of meromorphic basis}

Fix the original analytic coordinate $s_0$, and retain
\[
 \begin{gathered}
 O=\mathcal O_{X,s_0}=\mathbb C\{z\},\quad z=s-s_0,\quad
 F=\operatorname{Frac}O=O[1/z],\\
 \omega=O\,ds,\quad\widehat\omega=\mathbb C[[z]]\,ds,
 \quad C_\omega=\widehat\omega/\omega.
 \end{gathered}
 \tag{DC.6}
\]
Thus $F$ consists of convergent meromorphic germs, and is not the
formal Laurent-series field.  The injective inclusion
$\omega\hookrightarrow\widehat\omega$ is the Taylor-series map.
Its quotient is nonzero: $\sum_{n\ge0}n!z^n\,ds$ is a formal
series, while its terms fail to tend to zero at every fixed $z\ne0$.
It is therefore not an analytic germ.

Multiplication by $z$ is bijective on $C_\omega$.  To prove
surjectivity, subtract the constant coefficient of a formal
representative and divide the remaining formal series by $z$.
The removed constant lies in $\omega$.  For injectivity, if
$z\widehat f\,ds\in\omega$, the analytic germ on the left has
zero constant term and can be divided analytically by $z$;
the uniqueness of its Taylor series gives $\widehat f\,ds\in\omega$.
Every nonzero analytic germ is a unit times a power of $z$, so
$C_\omega$ is an $F$-vector space with the uniquely determined
action by inverses of these multiplications.

For every $F$-vector space $L$, viewed as an $O$-module, one has
\[
 \operatorname{Hom}_O(L,\widehat\omega)=0,
 \qquad \operatorname{Ext}^i_O(L,\widehat\omega)=0\quad(i\ge1).
 \tag{DC.7}
\]
The Hom assertion follows from
$f(l)\in\bigcap_nz^n\widehat\omega=0$, since $l=z^nl_n$.
For the Ext assertion first use the exact free resolution (OCQ.15)
of $F$.  Its Hom differential on the products of
$\widehat\omega$ is
$\Delta(a)_n=a_n-za_{n+1}$.  Given an arbitrary sequence
$(c_n)$ in $\widehat\omega$, set
\[
 a_n=\sum_{j\ge0}z^jc_{n+j}\in\widehat\omega.
 \tag{DC.8}
\]
For each fixed coefficient only finitely many summands contribute,
so this sum is well-defined in the specified completion.
Subtracting $za_{n+1}$ cancels every term except $c_n$.
Thus $\Delta$ is onto.  The resolution has length one, so all
higher Ext groups vanish.  A direct sum of these free resolutions
over an $F$-basis of $L$ is a free resolution of $L$; its Hom
differential is the product of the preceding surjections, hence
is onto.  This proves (DC.7) for arbitrary $L$.  The basis is used
only to prove this vanishing assertion.

An $O$-linear map between $F$-vector spaces is $F$-linear: applying
the map to $z(z^{-1}l)=l$ and using invertibility of $z$ in the
target proves compatibility with $z^{-1}$, and compatibility with
all meromorphic fractions follows.  Consequently the exact
completion sequence gives a canonical isomorphism, natural in $L$,
\[
 \begin{split}
 0\longrightarrow\omega\longrightarrow\widehat\omega
       \longrightarrow C_\omega\longrightarrow0,
 \qquad
 \partial_L:\operatorname{Hom}_F(L,C_\omega)
          \xrightarrow{\sim}\operatorname{Ext}^1_O(L,\omega).
 \end{split}
 \tag{DC.9}
\]
Indeed its Hom--Ext sequence has zero groups adjacent to this map
by (DC.7).  The map $\partial_L$ sends $f$ to the literal pullback
extension
\[
 E_f=\{(l,\widehat v)\in L\oplus\widehat\omega:
                      f(l)=[\widehat v]\},\qquad
 0\to\omega\xrightarrow{v\mapsto(0,v)}E_f
       \xrightarrow{(l,\widehat v)\mapsto l}L\to0.
 \tag{DC.10}
\]
Surjectivity of the last arrow follows from surjectivity onto
$C_\omega$, and its kernel is the displayed copy of $\omega$.
The extension is split exactly when $f$ lifts to an $O$-linear
map $L\to\widehat\omega$; by (DC.7), this occurs exactly when
$f=0$.  Thus the canonical description detects every such nonzero
map without choosing a basis.

For $L=K=\mathcal K_{s_0}$, evaluation at a chosen $F$-basis
identifies the left side of (DC.9) with
$\prod_B C_\omega$.  This recovers the original product calculation
(OCQ.21), with its full differential form factor.  The isomorphism
(DC.9) itself is independent of that choice.

\subsection{The natural local derived bidual retains every kernel class}

Define the local derived dual by
\[
 D_\omega L=R\operatorname{Hom}_O(L,\omega[1]).
 \tag{DC.11}
\]
Here complexes use cohomological grading; $\omega[1]$ has
$\omega$ in degree $-1$.  For an $F$-vector space $L$, its free
resolution has length one, its ordinary Hom into $\omega$ is zero,
and (DC.9) computes its Ext in degree one.  Therefore
\[
 D_\omega L\simeq\operatorname{Hom}_F(L,C_\omega)[0].
 \tag{DC.12}
\]
The module on the right is again an $F$-vector space, by pointwise
scalar multiplication.  Applying the same calculation a second
time yields
\[
 D_\omega^2 L\simeq
 \operatorname{Hom}_F(\operatorname{Hom}_F(L,C_\omega),C_\omega)[0].
 \tag{DC.13}
\]

Under these canonical identifications the natural derived bidual
map is exactly
\[
 \operatorname{ev}_L:L\longrightarrow D_\omega^2L,
 \qquad l\longmapsto\bigl(f\longmapsto f(l)\bigr).
 \tag{DC.14}
\]
To check the map and its sign, represent $\omega[1]$ by
\[
 J=[\widehat\omega\xrightarrow{-\pi}C_\omega],
 \qquad\deg\widehat\omega=-1,\quad\deg C_\omega=0.
 \tag{DC.15}
\]
The quasi-isomorphism $\omega[1]\to J$ is the positive Taylor
inclusion in degree $-1$.  Both terms of $J$ are Hom-acyclic for
an $F$-vector source: this is (DC.7) for $\widehat\omega$;
for $C_\omega$ its telescope differential is onto by the recursion
$a_0=0$, $a_{n+1}=z^{-1}(a_n-c_n)$, and the same direct-sum
resolution proves higher vanishing.  Thus
$\operatorname{Hom}_O(L,J)$ computes $D_\omega L$ and has only
the term $\operatorname{Hom}_F(L,C_\omega)$ in degree zero.
The identical reasoning applies to that $F$-vector space as source.
Evaluation into $J$ consequently computes the derived bidual map.
Its dg sign is $(-1)^{|l||f|}=1$, since both $l$ and $f$ have
degree zero.  It is the displayed positive evaluation.  Equivalently,
comparison of $J$ with an injective resolution gives
quasi-isomorphisms on these two Hom complexes by the proved
acyclicity, and carries this evaluation to the natural derived one.

The map (DC.14) is injective.  If $l\ne0$, extend $l$ to an
$F$-basis of $L$, and take its coordinate functional
$\lambda:L\to F$ with $\lambda(l)=1$.  Choose the nonzero class
$c=[\sum n!z^n\,ds]\in C_\omega$.  The $F$-linear map
$f(x)=\lambda(x)c$ satisfies $f(l)=c\ne0$, so
$\operatorname{ev}_L(l)\ne0$.  This proves injectivity of the
canonical map; the auxiliary basis and nonzero $c$ only verify it.
No assertion of surjectivity is made.

In particular every nonzero class in the original $K$ survives its
natural local derived bidual.  This includes the original class
$k_0=g\otimes[F_0]$ of (OCQ.5).  The original ordinary dual
vanishing $\operatorname{Hom}_O(M,\omega)=0$ therefore cannot be
used to discard this kernel from the derived comparison.
The finite summand is retained simultaneously.  Represent $N$ by
$P=[O\xrightarrow{g}O]$ in degrees $-1,0$.  Its dual is
$[\omega\xrightarrow{g}\omega]$ in the same degrees: the Hom
differential on a degree-minus-one map is $-(-1)^{-1}fg=fg$.
Dualizing again gives precisely $[O\xrightarrow{g}O]$.
In each term the evaluation sign is positive, since the possible
degree pairs are $(-1,0)$ and $(0,-1)$.  Finite free evaluation
is the identity in the displayed bases.  Consequently
$D_\omega N=\omega/g\omega$ in degree zero and
$N\xrightarrow{\sim}D_\omega^2N$.  Thus the canonical splitting
$M=N\oplus K$ makes its natural derived bidual injective, with
the entire finite arithmetic summand mapped isomorphically.

\subsection{The exact global target of the local Ext calculation}

Let $i_{s_0}$ denote the inclusion of the one-point ringed space with
ring $O$ into $(X,\mathcal O_X)$.  For an $O$-module $E$, its
direct image has sections $E$ on an open containing $s_0$ and zero
on an open not containing it, with the analytic action through germs.
This functor is exact.  The exact stalk functor is its left adjoint,
as one checks by restricting a sheaf map to germs, so the direct
image preserves injective modules.  Applying it to an injective
resolution proves the exact derived adjunction
\[
 R\operatorname{Hom}_{\mathcal O_X}(\mathcal M,i_{s_0*}\omega)
 \simeq R\operatorname{Hom}_O(M,\omega).
 \tag{DC.16}
\]
The same proof on every open gives the internal version with
$i_{s_0*}$ on the right.  This places the explicit local extension
(DC.10) in a genuine sheaf-derived target.  That target
$i_{s_0*}\omega$ is not coherent on $X$: a coherent sheaf supported
at a point has a stalk killed by a power of $z$, whereas no such
power kills $\omega$.  Formula (DC.16) does not replace the
coherent target $\omega_X$ by a claimed equality.  The genuine
coherent sheaf target retains the globally split summand (DC.5),
independently of this skyscraper realization.

\subsection{The original tau complex, its dilation, and its retained degree zero}

The source in A1523 has the two incoming amplitude legs $V$,
the central amplitude $\mathscr B$, restrictions $\Theta$ and
$\Theta\widehat{\phantom\phi}$, and its cohomology complex
$[V\oplus V\to\mathscr B]$ with differential
$(\phi,\psi)\mapsto\Theta(\phi-\widehat\psi)$.
For its balanced action give the plus leg $V_+$ the action
$t=D$, and the minus leg $V_-$ the action $t=1-D$.
With the original Fourier kernel $e^{-2\pi ix\nu}$,
$\widehat{D\psi}=(1-D)\widehat\psi$ and Fourier squares to the
identity on these even functions.  Therefore Fourier is the exact
$A$-linear isomorphism $V_-\to V_+$.

Write $\mathcal V_+$, $\mathcal V_-$, and $\mathcal B$ for the
sheafifications of the corresponding $\mathcal O_X\otimes_A$
presheaves.  Flatness of every analytic stalk over $A$ proves that
\[
 C_A=[\mathcal V_+\xrightarrow{\Theta}\mathcal B],\qquad
 C_\tau=[\mathcal V_+\oplus\mathcal V_-
       \xrightarrow{(\phi,\psi)\mapsto\Theta(\phi-\widehat\psi)}
        \mathcal B]
 \tag{DC.17}
\]
have their indicated original differentials, with terms in degrees
zero and one.  There is an explicit cochain isomorphism
\[
 C_\tau\xrightarrow{\sim}C_A\oplus\mathcal V_+[0],
 \qquad (\phi,\psi)\longmapsto
             (\phi-\widehat\psi,\widehat\psi)
 \tag{DC.18}
\]
in degree zero and the identity on $\mathcal B$ in degree one.
Its inverse sends $(u,v)$ to $(u+v,\widehat v)$.  Applying the
two differentials gives $\Theta u$ in both directions, proving
the cochain assertion.  The quotient map in degree one is a
quasi-isomorphism $C_A\to\mathcal M[-1]$, because the original
$\Theta$ is injective and balancing is flat.  Thus (DC.3) gives
the canonical derived decomposition
\[
 C_\tau\simeq\mathcal V_+[0]\oplus\mathcal N[-1]
                         \oplus\mathcal K[-1].
 \tag{DC.19}
\]
The degree-zero diagonal source is retained.  On the actual
degree-one tau cohomology, (DC.19) is exactly the split comparison
$\mathcal M=\mathcal N\oplus\mathcal K$.

The comparison to the analytic complex
$C_g=[\mathcal O_X\xrightarrow{g}\mathcal O_X]$ is also an
actual cochain map: on $C_A$ it sends $a\otimes\phi$ to
$aH_\phi$ and $a\otimes F$ to $a\mathcal MF$.
The original identity $\mathcal M\Theta=gH$ proves that it
commutes with the differentials.  On $C_\tau$ compose it with
the first projection in (DC.18).  The diagonal source is exactly
the degree-zero kernel of that first projection.  The full composite
has the larger kernel computed next; no injectivity of $H$ is inferred.
Its induced degree-one map is precisely $\beta$.

Write $H:\mathcal V_+\to\mathcal O_X$ and
$\mathrm{Mel}:\mathcal B\to\mathcal O_X$ for the two displayed
maps, and put
\[
 A_0=\ker H,\qquad A_1=\ker\mathrm{Mel}.
 \tag{DC.20}
\]
Both maps are onto as sheaf maps.  For $H$ use the original
$\phi_*=(4\pi^2x^4-6\pi x^2)e^{-\pi x^2}$ with $H_{\phi_*}=1$;
the section $a$ lifts to $a\otimes\phi_*$.  For $\mathrm{Mel}$
use $aH_0^{-1}\otimes F_0$ with the original nowhere-zero
$H_0=\sqrt\pi e^{s^2/4}$.  In (DC.18) coordinates the full
kernel complex and its actual short exact sequence are therefore
\[
 C_{\ker}=[A_0\oplus\mathcal V_+
       \xrightarrow{(u,v)\mapsto\Theta u}A_1],\qquad
 0\to C_{\ker}\to C_\tau\to C_g\to0.
 \tag{DC.21}
\]
The inclusion $\Theta A_0\subset A_1$ follows from
$\mathrm{Mel}\Theta=gH$.  Since $\Theta$ is injective,
$H^0C_{\ker}=\mathcal V_+$, its diagonal factor.  The map
$A_1\to\mathcal K$ sends a target vector to its original
theta-quotient class.  It is onto locally: if a representative $F$
has $\mathrm{Mel}(F)=ga$, subtract
$\Theta(a\otimes\phi_*)$ to obtain a representative in $A_1$.
If a vector in $A_1$ is $\Theta\phi$, then $gH(\phi)=0$;
injectivity of $g$ on $\mathcal O_X$ gives $\phi\in A_0$.
Thus its kernel is exactly $\Theta A_0$.  This proves
\[
 H^0C_{\ker}=\mathcal V_+,\qquad
 H^1C_{\ker}=A_1/\Theta A_0\xrightarrow{\sim}\mathcal K,
 \qquad
 C_{\ker}\simeq\mathcal V_+[0]\oplus\mathcal K[-1].
 \tag{DC.22}
\]
For the last assertion project the explicit complex onto its
diagonal factor and take the quotient of its remaining degree-one
term.  The other degree-zero differential is injective, so the
resulting map is a quasi-isomorphism.

Finally the original dilation $R_aF(x)=F(x/a)$ acts as $R_a$ on
$\mathscr B$ and $V_+$ and as $aR_{1/a}$ on $V_-$.
Direct Fourier substitution gives
$\widehat{aR_{1/a}\psi}=R_a\widehat\psi$, so (DC.18) is
equivariant with these exact actions.  The Mellin substitution
$x=ay$ gives $\mathcal M R_aF=a^s\mathcal MF$ and
$H_{R_a\phi}=a^sH_\phi$.  Hence the cochain map to $C_g$
intertwines these actions with multiplication by $a^s$ on both
analytic terms.  Dilation commutes with $g$ and preserves
$\mathcal K$, so it commutes with the projections (DC.2);
in particular $R_a j=j a^s$ on $\mathcal N$.

At an actual zero $\rho$ with its full order $m$ and full local
unit $g(s)=u(s-\rho)(s-\rho)^m$, this action is exactly
\[
 a^\rho\sum_{j=0}^{m-1}\frac{(\log a)^j}{j!}N_\rho^j,
 \qquad N_\rho[v]=(s-\rho)[v]\quad\text{in }O/(g).
 \tag{DC.23}
\]
No nilpotent term is removed.  The finite coherent dual summand
(DC.5) passes by the original maps
$[v\,ds]\mapsto[v\,ds/g]$ and residue to the original tau
dual target exactly as proved in (OCQ.26)--(OCQ.30).  Its dual
dilation is $\lambda\mapsto(q\mapsto a\lambda(R_{1/a}q))$;
the differential representative is multiplied by $a^{1-s}$,
including the retained one-dimensional character.

For a proper admitted source $W$, the entire relation quotient
$\mathcal O_X\otimes_A(V/W)$ and its source-label transition
remain the ones in (OCQ.31)--(OCQ.33).  No invertibility on that
larger proper-label kernel is used in (DC.1)--(DC.23).
The diagonal degree-zero summand in (DC.19) belongs to the joint
face $S=\{+,-\}$.  Under (DC.18), the plus-only source is exactly
the subspace $\{(u,0)\}$, and the minus-only source is exactly
$\{(-v,v)\}$.  On these subspaces the differential is respectively
$\Theta u$ and $-\Theta v$, so each one-leg complex has
$H^0=0$ and $H^1=\mathcal M$.  These are the specified restrictions
of the cochain maps; no joint-face diagonal is added to a one-leg
face.  On each such original support fibre the restricted linear
maps lift by $(S,v)\mapsto(S,f(v))$ and external
$\tau\mapsto\tau$, as in (OCQ.34).  Thus a killed vector lands
at its supported zero.
These exact cochain, coherent, residue, and derived-dual maps
retain the actual counterfactual packet throughout; they do not
supply a new weight bound or assert the existence of such a zero.
